\documentclass[10pt]{article}
\usepackage[T1]{fontenc}
\usepackage[utf8]{inputenc}
\usepackage[french]{babel}
\usepackage{amsmath,amssymb,mathrsfs,mathpazo}
\usepackage{geometry}
\geometry{a4paper,margin=2cm}
\usepackage{graphicx,array,tabularx,tabularray,colortbl,multirow,multicol,booktabs,enumitem,pifont,ulem,fancybox,diagbox,cancel}
\usepackage{pstricks,pstricks-add,pst-plot,pst-tree,pst-node,pst-text,pst-func,pst-math,pst-bspline,pst-3dplot,pst-eucl,pst-solides3d}
\let\pscalculate\undefined
\usepackage{pst-bezier}
\usepackage{tikz,pgfplots,tkz-tab}
\usetikzlibrary{arrows,arrows.meta,patterns,calc,positioning,trees,decorations.markings,shapes.geometric,shapes.misc,shapes.arrows,intersections,decorations.pathmorphing,decorations.pathreplacing,backgrounds,shadows,babel}
\pgfplotsset{compat=1.18}
\usepgfplotslibrary{fillbetween}
\usepackage[np]{numprint}
\usepackage{listings,makecell,pas-tableur,twoopt}
\let\barre\undefined
\usepackage[pstricks]{tablvar}
\definecolor{darkgreen}{RGB}{0,100,0}
\providecommand{\red}{\color{red}}
\providecommand{\blue}{\color{blue}}
\providecommand{\green}{\color{green}}
\providecommand{\black}{\color{black}}
\providecommand{\white}{\color{white}}
\providecommand{\gray}{\color{gray}}
\providecommand{\cyan}{\color{cyan}}
\providecommand{\magenta}{\color{magenta}}
\providecommand{\yellow}{\color{yellow}}
\providecommand{\violet}{\color{violet}}
\providecommand{\orange}{\color{orange}}
\providecommand{\euro}{\texteuro}
\providecommand{\EUR}{\texteuro}
\providecommand{\ieme}{\textsuperscript{e}}
\providecommand{\degres}{\ensuremath{^\circ}}
\providecommand{\degree}{\ensuremath{^\circ}}
\providecommand{\vv}[1]{\overrightarrow{#1}}
\providecommand{\Coord}[1]{#1}
\providecommand{\decosix}{}
\providecommand{\starredbullet}{\ensuremath{\bullet}}
\pagestyle{empty}
\setlength{\parindent}{0pt}
\usepackage{adjustbox}\DeclareUnicodeCharacter{2236}{:}\DeclareUnicodeCharacter{2190}{\ensuremath{\leftarrow}}
\providecommand{\R}{\mathbb{R}}
\providecommand{\N}{\mathbb{N}}
\providecommand{\Z}{\mathbb{Z}}
\providecommand{\Q}{\mathbb{Q}}
\providecommand{\D}{\mathbb{D}}
\providecommand{\C}{\mathbb{C}}
\providecommand{\e}{{\mathrm{e}}}
\providecommand{\vect}[1]{\overrightarrow{#1}}
\providecommand{\vectt}[1]{\overrightarrow{\mathrm{#1}}}
\providecommand{\vvt}[1]{\overrightarrow{\mathrm{#1}}}
\providecommand{\np}[2][]{{#2~#1}}
\providecommand{\vv}[1]{{\overrightarrow{#1}}}
\renewcommand{\d}{{\mathrm{d}}}
\providecommand{\degres}{{^\circ}}
\providecommand{\degre}{{^\circ}}
\providecommand{\blue}{\color{blue}}
\providecommand{\red}{\color{red}}
\providecommand{\enskip}{\quad}
\providecommand{\llbracket}{\lbrack\!\lbrack}
\providecommand{\rrbracket}{\rbrack\!\rbrack}
\providecommand{\euro}{{\text{€}}}
\providecommand{\no}{{\text{n°}}}
\providecommand{\interval}[2]{{[#1;#2]}}
\providecommand{\intervFO}[2]{{[#1;#2[}}
\providecommand{\dotfill}{\ldots}

\providecommand{\ds}{\displaystyle}
\providecommand{\textding}[1]{#1}
\providecommand{\Oij}{\(\mathrm{O};\vec{\imath},\vec{\jmath}\)}
\providecommand{\Oijk}{\(\mathrm{O};\vec{\imath},\vec{\jmath},\vec{k}\)}
\providecommand{\Ouv}{\(\mathrm{O};\vec{u},\vec{v}\)}
\providecommand{\intervFF}[2]{\ensuremath{[#1;#2]}}
\renewcommand{\euro}{\ifmmode\text{\texteuro}\else\texteuro\fi}
\DeclareUnicodeCharacter{00B0}{\ifmmode^{\circ}\else\textdegree\fi}
\DeclareUnicodeCharacter{20AC}{\ifmmode\text{\texteuro}\else\texteuro\fi}
% Énoncés : Éducation nationale. Transcription et corrigés : APMEP, Denis Vergès et contributeurs. Découpage et classement : CoopMaths, Sébastien Lozano, Éric Elter et contributeurs. Adaptation de la présentation : Maths Explorer.
% CC BY-SA : https://coopmaths.fr/www/about/
% https://coopmaths.fr/alea/static/bac/2023/tex/bac_2023_03_sujet1_reunion_2_cor.tex
\begin{document}
\noindent\textit{APMEP / CoopMaths. CC BY-SA. Présentation adaptée par Maths Explorer.}\par\medskip
\medskip

On considère la fonction $f$ définie sur $]0~;~+\infty[$ par: $f(x) = 3x + 1 - 2x \ln (x)$.

On admet que la fonction $f$ est deux fois dérivable sur $]0~;~+\infty[$. 
On note $f'$ sa dérivée et $f''$ sa dérivée seconde.
On note $\mathcal{C}_f$ sa courbe représentative dans un repère du plan.



\begin{enumerate}
\item On cherche les limites de la fonction $f$ en 0 et en $+\infty$.

\begin{list}{\textbullet}{}
\item On sait que $\ds\lim_{x\to 0 \atop x>0} x\ln(x)=0$ donc $\ds\lim_{x\to 0 \atop x>0} \left (3x+1 -2x\ln(x)\right )=1$ et donc $\ds\lim_{x\to 0 \atop x>0} f(x)=1$.
\item $f(x)=x\left ( 3-2\ln(x)\right ) +1$.

$\ds\lim_{x\to +\infty} \ln(x)=+\infty$ donc $\ds\lim_{x\to +\infty} \left (3-2\ln(x)\right )=-\infty$.

On en déduit que $\ds\lim_{x\to +\infty} x\left (3-2\ln(x)\right )=-\infty$ et donc que $\ds\lim_{x\to +\infty} f(x)=-\infty$.
\end{list}
\item 
	\begin{enumerate}
		\item 
Sur $]0~;~+\infty[$, $f'(x)=	3+0-2\times\ln(x) -2x\times \dfrac{1}{x}=3-2\ln(x)-2 = 1-2\ln(x)$.
		
		\item On étudier le signe de $f'$  sur l'intervalle $]0~;~+\infty[$.
		


$f'(x)>0 \iff 1-2\ln(x)>0 \iff \dfrac{1}{2}>\ln(x) \iff x<\e^{\frac{1}{2}}$

$f\left (\e^{\frac{1}{2}}\right ) = 3 \e^{\frac{1}{2}} +1 -2 \e^{\frac{1}{2}} \times \ln\left (\e^{\frac{1}{2}}\right ) = 3 \e^{\frac{1}{2}} +1 -2 \e^{\frac{1}{2}} \times \dfrac{1}{2} =
3 \e^{\frac{1}{2}} +1 -  \e^{\frac{1}{2}} = 2  \e^{\frac{1}{2}} +1 \approx 4,3$.

On en déduit le tableau des variations de $f$ sur l'intervalle $]0~;~+\infty[$.

\begin{center}
{\renewcommand{\arraystretch}{1.3}
\psset{nodesep=3pt,arrowsize=2pt 3}
\def\esp{\hspace*{0.8cm}}
\def\hauteur{20pt}
$\begin{array}{|c|l*6{c}|}
\hline
x & 0  & \esp \esp & \e^{\frac{1}{2}} & \esp &  & \esp & +\infty \\ 
\hline
\text{ Signe de $f'$ } & \vline\;\vline  &   \pmb{+}   &   \vline\hspace{-2.7pt}0 &  & \pmb{-} &  &\\ 
\hline
 &\vline\;\vline  & & \Rnode{max}{2\e^{\frac{1}{2}}+1} &&  &   &     \\  
\text{ Variations de $f$ } & \vline\;\vline  && & &  &  &  \rule{0pt}{\hauteur} \\ 
 & \vline\;\vline  ~\Rnode{min1}{1}&   &  &  &   & &  \Rnode{min2}{-\infty}  \rule{0pt}{\hauteur}    
 \ncline{->}{min1}{max}   \ncline{->}{max}{min2}
 \\ 
\hline
\end{array} $
}
\end{center}
	\end{enumerate}	

\item
	\begin{enumerate}
		\item 
Le maximum de la fonction $f$ vaut $2 \e^{\frac{1}{2}} +1>0$; on complète le tableau des variations de $f$.
		
\begin{center}
{\renewcommand{\arraystretch}{1.3}
\psset{nodesep=3pt,arrowsize=2pt 3}
\def\esp{\hspace*{0.8cm}}
\def\hauteur{20pt}
$\begin{array}{|c|l*6{c}|}
\hline
x & 0  & \esp \esp & \e^{\frac{1}{2}} & \esp & \Rnode{alpha}{\blue\alpha} & \esp & +\infty \\ 
\hline
 &\vline\;\vline  & & \Rnode{max}{2\e^{\frac{1}{2}}+1} &&  &     &   \\  
\text{ Variations de $f$ } & \vline\;\vline  && & &\Rnode{zero1}{\blue{}0}   &  &  \rule{0pt}{\hauteur} \\ 
 & \vline\;\vline  ~\Rnode{min1}{1}&   &  &  &   & &  \Rnode{min2}{-\infty}  \rule{0pt}{\hauteur}    
 \ncline{->}{min1}{max} 
  \ncline{max}{zero1}  \ncline{->}{zero1}{min2}
  \ncline[linecolor=blue,linestyle=dashed]{zero1}{alpha}
 \\ 
\hline
\end{array} $
}
\end{center}

La fonction $f$ est dérivable sur $\left ] 0\;;\;+\infty\right [$ donc continue sur cet intervalle.
		
\begin{list}{\textbullet}{}
\item Sur l'intervalle $\left ] -\infty~,~\e^{\frac{1}{2}}\right [$, la fonction $f$ est strictement positive donc l'équation $f(x)=0$ n'a pas de solution.
\item Sur l'intervalle $\left ] \e^{\frac{1}{2}}\;;\;+\infty\right [$, la fonction $f$ est continue et strictement décroissante; elle passe d'une valeur positive à une valeur négative donc, d'après le corollaire du théorème des valeurs intermédiaires, l'équation $f(x)=0$ admet une solution unique sur cet intervalle. On l'appelle $\alpha$.
\end{list}

On en déduit que l'équation $f(x)=0$ admet une solution unique $\alpha$ sur $]0\;;\;+\infty[$, et que cette solution appartient à l'intervalle $\left ] \e^{\frac{1}{2}}\;;\;+\infty\right [$.
		
		\item Le signe de la fonction $f$ sur $]0~;~+\infty[$ est donné par le tableau:
		
\begin{center}
{\renewcommand{\arraystretch}{1.3}
\psset{nodesep=3pt,arrowsize=2pt 3}
\def\esp{\hspace*{0.8cm}}
\def\hauteur{20pt}
$\begin{array}{|c|l*6{c}|}
\hline
x & 0  & \esp \esp & \e^{\frac{1}{2}} & \esp & \Rnode{alpha}{\blue\alpha} & \esp & +\infty \\ 
\hline
 &\vline\;\vline  & & \Rnode{max}{2\e^{\frac{1}{2}}+1} &&  &     &   \\  
\text{ Variations de $f$ } & \vline\;\vline  && & &\Rnode{zero1}{\blue{}0}   &  &  \rule{0pt}{\hauteur} \\ 
 & \vline\;\vline  ~\Rnode{min1}{1}&   &  &  &   & &  \Rnode{min2}{-\infty}  \rule{0pt}{\hauteur}    
 \ncline{->}{min1}{max} 
  \ncline{max}{zero1}  \ncline{->}{zero1}{min2}
  \ncline[linecolor=blue,linestyle=dashed]{zero1}{alpha}
 \\ 
 \hline
\text{ Signe de $f$ } & \vline\;\vline &  &   \pmb{+} & & \vline\hspace{-2.7pt}0 &  \pmb{-}  &\\ 
\hline
\end{array} $
}
\end{center}

Donc $f(x)>0$ sur $]0\;;\;\alpha[$, $f(\alpha)=0$ et $f(x)<0$ sur $]\alpha\;;\;+\infty[$.
		
		
	\end{enumerate}	
\item On considère une primitive quelconque de la fonction $f$ sur
l'intervalle $]0~;~+\infty[$. On la note $F$.



La fonction $F$ a pour dérivée la fonction $f$ qui est positive sur $\left[\e^{\frac12}~;~\alpha\right[$; donc $F$ est croissante sur cet intervalle, ce qui contredit l'affirmation proposée.
\item
	\begin{enumerate}
		\item Pour étudier la convexité de la fonction $f$ sur $]0~;~+\infty[$, on calcule la dérivée seconde $f''$:
		
$f''(x)=0-2\times \dfrac{1}{x}=-\dfrac{2}{x}<0$ sur $]0~;~+\infty[$; on en déduit que la fonction $f$ est concave sur $]0~;~+\infty[$ et que sur cet intervalle,  la courbe $\mathcal{C}_f$ est située au dessous des ses tangentes.
		\item Soit $T$ la tangente à la courbe $\mathcal{C}_f$ au point d'abscisse 1.

$T$ a pour équation $y=f'(1)\left (x-1\right ) + f(1)$.

~~\textbullet~~$f(x)=3x+1-2x\ln(x)$ donc $f(1)=3\times 1 +1 -2\times 1 \times \ln(1)=4$

~~\textbullet~~$f'(x)=1-2\ln(x)$ donc $f'(1)=1-2\times\ln(1)=1$

Donc $T$ a pour équation $y=1\left (x-1\right ) +4$ soit $y=x+3$.
		
		\item 
La courbe $\mathcal{C}$ est en dessous de ses tangentes donc de la tangente $T$ et donc, pour tout $x$ de $]0\;;\;+\infty[$, on a $f(x)\leqslant x+3$.

$\begin{gathered}f(x)\leqslant x+3
\iff 3x+1-2x\ln(x) \leqslant x+3
\iff 2x-2 \leqslant 2x\ln(x)
\iff \dfrac{2x}{2x} - \dfrac{2}{2x} \leqslant \ln(x)\\
\phantom{f(x)\leqslant x+3}
\iff 1-\dfrac{1}{x} \leqslant \ln(x)\end{gathered}$

Donc, pour tout  $x$ de $]0\;;\;+\infty[$, on a: $\ln(x)\geqslant 1-\dfrac{1}{x}$.
		
	\end{enumerate}
\end{enumerate}

\bigskip

\end{document}
